\chapter{Những Quy Định Trên Trời}
\markboth{Những Quy Định Trên Trời}{Những Quy Định Trên Trời}

\begin{truyen}{Tóc Dài Có Làm Sai Toán Không?}{Questioning School Rules}
\chuhoa{M}{ột bạn nam bị nhắc vì tóc dài, liền hỏi ngay trước lớp: ``Cô cho em hỏi thật, tóc em dài hơn quy định thì nó làm em giải sai Toán ở chỗ nào ạ?''}

\teacherpair{Cô không né: ``Nó không làm em sai Toán trực tiếp. Nhưng trường học không chỉ dạy kết quả bài kiểm tra. Nó còn dạy cách em sống trong một tập thể có khuôn khổ. Vấn đề là nhiều khi người lớn giải thích nội quy quá tệ nên các em chỉ thấy phần cấm mà không thấy phần lý do.''}{The teacher answers: ``It does not directly make you solve math wrongly. But school is not only about test results. It also teaches how to live in a community with shared structure. The problem is that adults often explain rules poorly, so you only see the prohibition and not the reason.''}

Bạn ấy vẫn chưa chịu: ``Vậy nếu lý do không thuyết phục thì tụi em có quyền hỏi chứ ạ?''

\teacherpair{Có quyền hỏi. Và người lớn nên tập trả lời cho tử tế hơn. Nhưng hỏi để hiểu khác với hỏi chỉ để phá luật cho sướng miệng. Nếu em muốn phản biện, hãy phản biện bằng thái độ khiến người khác muốn nghe em đến cùng.}{``You do have the right to ask. Adults should learn to answer more properly. But asking to understand is different from asking just to mock the rule. If you want to challenge something, do it with a manner that makes people willing to hear you out.''}

Cô ngừng một chút rồi nói thật hơn: ``Nói cho công bằng, có những nội quy chính thầy cô cũng không phải người đặt ra. Có khi tụi cô đang đứng giữa: một bên là quy định từ trên xuống, một bên là câu hỏi rất có lý của các em. Vậy nên càng không thể lấy thái độ cáu gắt để bù cho một lời giải thích yếu.''

\textit{(A student questions a rule. The teacher admits the rule may be poorly explained, while still defending the need for structure.)}
\end{truyen}

\begin{baihoc}
Nội quy muốn có sức nặng thì phải đi cùng lý do rõ ràng. Phản biện muốn có giá trị thì phải đi cùng thái độ đàng hoàng. Và người lớn muốn được tôn trọng thì cũng phải dám giải thích điều mình đang buộc người khác làm theo.
\textit{(Rules carry weight only when paired with clear reasons. Critique carries value only when paired with proper attitude. Adults who want respect must also explain what they ask others to obey.)}
\end{baihoc}

\begin{ghinhoanh}
\textbf{Short English to remember:}

Rules need reasons.

Questions need respect.
\end{ghinhoanh}

\begin{tuhocnhanh}
1. Em đang phản biện để hiểu hơn hay để thắng miệng? \textit{(Are you questioning to understand, or just to win verbally?)}

2. Người lớn quanh em có nội quy nào nên được giải thích lại cho hợp lý hơn? \textit{(Which rule around you should be explained more reasonably?)}
\end{tuhocnhanh}

\ngancach